\section{Robustness of the edit obstruction under marker permutations}
\label{sec:robust}
Only a small part of the complete marker order enters the edit-distance
argument.  The original host places the two groups of each marker
mode consecutively and uses a particular ordering of all $6h$ modes.
The forcing argument itself requires only that the two groups of a
single mode occur in their given order and that the entire marker
prefix precede the variable and dummy blocks. This gives a useful
separation between the mechanism responsible for distance and the
mechanism responsible for terminal localization.

For each depth $h$, let $\pi_r$ and $\pi_c$ be two arbitrary
permutations of its $6h$ complete marker modes. Define
$A_h^{\pi_r,\pi_c}$ by permuting the marker \emph{row} blocks in
accordance with $\pi_r$ and marker \emph{column} blocks in
accordance with $\pi_c$, preserving the order of the two groups
inside each mode. Keep every block-constant entry, the two signed
tree orders, and the positions of the row and column dummy blocks
relative to the variable tree. The result is again a square matrix
of order $(14h+2)2^h$.

\begin{proposition}[Order-independent forcing lower bound]
\label{prop:permuted}
For every $h\ge1$ and every pair of complete-mode permutations,
\[
 D_{H_4}(A_h^{\pi_r,\pi_c})\ge 4^h.
\]
In particular, the lower bound holds against arbitrary binary edits
under independently reordered marker prefixes.
\end{proposition}
\begin{proof}
Use exactly the $m^2$ XOR-indexed selections from
Section~\ref{sec:distance}. The definition and injectivity of the
protected cells depend on physical block labels, their offsets,
the chosen tree path and the fact that marker and dummy blocks have
size $m$; permutation of whole marker modes changes none of these.
Therefore any edit set of size below $m^2$ misses the protected cells
of one selection.

For a fixed mode $t$, its two marker row groups appear in the order
$A_1(t)<A_2(t)$ and its marker column groups in the order
$B_1(t)<B_2(t)$. All four groups precede every variable and dummy
position, regardless of the position of other complete modes in the
prefix. The four forcing frames use only the two marker groups of
their own mode. Their increasing order and all the relevant
incidences therefore remain unchanged. The implications in the
forcing table again propagate $x_0=1$ and $y_0=0$ to the common leaf
cell, contradicting freeness. Thus any eliminating edit set has
size at least $m^2$.
\end{proof}

The upper bound $D_{H_4}(A_h)\le m^2$ and the exact copy count were
proved for the canonical mode order in Section~\ref{sec:host}.
Proposition~\ref{prop:permuted} records the broader range of the
forcing mechanism itself. The separately supplied Lean extension
contains the corresponding concrete-index theorem
\texttt{permuted\_matrix\_edit\_lower\_bound}; the primary
submission project retains its canonical order in the global
classification and removal statements.

\section{The selection map as an injective coding device}
\label{sec:selectionmap}
The XOR choice can be characterized without mentioning the particular
marker rules. Fix a word length $h$, and for each pair
$(a,b)\in\{0,\dots,2^h-1\}^2$ place a selected representative in
every signed variable block on the ancestral path of $a\oplus b$.
At a depth $d$, the selected node is the high-$d$-bit prefix of the
XOR, and the selected offset on the row axis (respectively column
axis) is the low $h-d$ bits of $a$ (respectively $b$).

\begin{lemma}[One-endpoint recovery]\label{lem:recovery}
Suppose the selected variable row at depth $d$ is known and the
selected marker-column offset is known. Then the pair $(a,b)$ is
uniquely determined. The analogous assertion holds with the two
axes reversed. A pair of selected root-variable positions also
uniquely determines $(a,b)$.
\end{lemma}
\begin{proof}
The marker-column offset equals $b$. The variable row determines
its node prefix $p_d$ and row offset $u=a\bmod2^{h-d}$. Since
bitwise XOR commutes with taking high-bit prefixes,
\[
 \lfloor a/2^{h-d}\rfloor
     =p_d\oplus\lfloor b/2^{h-d}\rfloor.
\]
The quotient on the left and remainder $u$ recover $a$. This
also holds for $d=h$, when the row offset is zero. Reversing the
axes recovers $b$ by the same equation. At $d=0$, the root node
is zero and the row and column offsets are exactly $a$ and $b$.
\end{proof}

This coding interpretation makes the protected-cell injectivity
transparent. Each protected cell contains at least one such
recoverable configuration; hence it belongs to at most one selected
path. The proof uses all $m^2$ labels simultaneously. Independent
choices of a random leaf and unrelated row and column offsets would
not provide this exact disjointness, even though some marginal
probabilities would still be uniform. The deterministic injection
also eliminates the need to estimate correlations among the
successive forcing gadgets.
